% TOP_PROBLEM: {"problemId": "problem.furstenberg-times-two-times-three-invariant-measures", "canonicalTitle": "Furstenberg's x2 x3 conjecture on invariant measures", "releaseRank": 183, "primaryDomain": "probability_ergodic_dynamics", "primaryDomainLabel": "Probability, ergodic theory & dynamics", "displayStatus": "open", "releaseStatus": "open"}
% !TeX program = lualatex
% Definition and sources only; this file is not an LLM solution attempt.
\documentclass[11pt]{article}
\usepackage[margin=25mm]{geometry}
\usepackage{fontspec}
\setmainfont{Latin Modern Roman}
\usepackage{amsmath,amssymb,mathtools}
\usepackage{unicode-math}
\setmathfont{Latin Modern Math}
\usepackage{xurl,hyperref}
\hypersetup{colorlinks=true,urlcolor=blue}
\setcounter{secnumdepth}{0}
\setlength{\parindent}{0pt}
\setlength{\parskip}{5pt}
\setlength{\emergencystretch}{3em}
\begin{document}
\section*{\texorpdfstring{Furstenberg's x2 x3 conjecture on invariant measures}{Furstenberg's x2 x3 conjecture on invariant measures}}\label{183}
\subsection{Definitions and mathematical statement}
On $\mathbb T={\mathbb{R}}/{\mathbb{Z}}$, let $T_2(x)=2x$ and $T_3(x)=3x$. A Borel probability measure $\mu$ is invariant if $(T_2)_*\mu=(T_3)_*\mu=\mu$. Joint ergodicity means every measurable set invariant modulo null sets under both maps has measure zero or one. The selected assertion is
\[
 \mu\text{ invariant and jointly ergodic}
 \quad\Longrightarrow\quad
 \mu=\text{Haar measure or }|\operatorname{supp}\mu|<\infty.
\]
Ergodicity for the joint action must not be silently strengthened to ergodicity for each map separately. The catalogue also mentions a coprime $(a,b)$ formulation; that broader parameter statement is recorded in its quotation, but an equivalence with the single pair $(2,3)$ is not established merely by these definitions.
\par\smallskip\textit{Formulation and concept references:} \href{https://ar5iv.labs.arxiv.org/html/2410.22701}{[S1]}.

\subsection{Short English statement}
Must a jointly ergodic circle measure preserved by both doubling and tripling be either uniform or concentrated on finitely many points?
\subsection{Sources}
\begin{itemize}
\item[S1]P. Burton, J. Panangaden, 'Formulations of Furstenberg's x2 x3 conjecture in complex analysis and operator algebras', arXiv:2410.22701 (2024).
\url{https://ar5iv.labs.arxiv.org/html/2410.22701}.
\textit{Formulation source.}
\end{itemize}
\end{document}
